% ---------------------------------------------------------------
\section{Related Work}
\label{sec:related}

\subsection{Unsupervised Drift Detection}
Two-sample tests are the standard label-free monitor. Univariate KS tests with a
multiple-testing correction remain strong baselines for tabular shift
detection~\cite{rabanser2019failing}. Kernel tests such as MMD target multivariate
differences~\cite{gretton2012kernel}. ADWIN maintains an adaptive window over a
stream of real values, such as errors, and detects changes in their mean~\cite{bifet2007adwin}.
ADWIN-U extends adaptive windowing to the unsupervised setting~\cite{assis2025adwinu}.
DriftLens models deep embeddings per label with Gaussians and measures their Fréchet
distance from a reference window, which gives real-time detection and drift curves on
unstructured data~\cite{greco2025driftlens,greco2024driftlensdemo}. Recent benchmarks
show that unsupervised detectors behave very differently on real streams, where onset
truth is rarely available~\cite{lukats2025benchmark,souza2020challenges,cerqueira2026framework}.
Gower-Winter et al.\ argue that window-based detection is ill-posed without
assumptions~\cite{gowerwinter2026window}. Uncertainty-based, label-free and conformal detectors continue to be
proposed~\cite{lu2025pudd,tan2025flexible,eliades2025conformal}.
\emph{Strength:} they need no labels, scale to many features and come with
calibrated false-alarm control. \emph{Limitation for this setting:} they report
alarms, and whether an alarm coincides with lost accuracy lies outside their
evaluation target.

\subsection{Drift Localization and Explanation}
Hinder et al.\ survey how drift can be located and explained and give model-based
explanations of drift~\cite{hinder2024surveyA,hinder2024surveyB,hinder2023modelbased}.
Regional drift detection localizes change in feature space~\cite{liu2017regional}.
The domino drift effect (DDE) infers drift in unmonitored features from correlated,
already-drifting monitored features, which saves per-feature detectors in high
dimension~\cite{szucs2025dde}. Causal explanations of drift aim at actionable
attribution~\cite{komnick2026causal}. In medical imaging, shift identification is used
to support root-cause analysis~\cite{roschewitz2026shiftid}.
\emph{Strength:} they turn an alarm into an inspectable hypothesis and, in the case of
DDE, cut monitoring cost in high dimension. \emph{Limitation:} explanations describe
distributional change. They are seldom tested on whether the explained change caused
the loss of performance, and they cannot address changes that leave $P(X)$ untouched.

\subsection{Performance Estimation without Labels}
Average confidence, the ATC threshold estimator and difference-of-confidence estimate
accuracy under shift from unlabelled data~\cite{garg2022atc,guillory2021predicting}.
Agreement-on-the-line estimates accuracy from the agreement between
models~\cite{baek2022agreement}. Kivimäki et al.\ study confidence-based estimators for
model monitoring~\cite{kivimaki2025confidence}. Label-shift correction assumes a fixed
$P(X\mid Y)$~\cite{lipton2018detecting}. All of these depend on calibration, which
often degrades under shift~\cite{guo2017calibration,zadrozny2002transforming,platt2000probabilities}.
Performance-aware detectors~\cite{bayram2022concept} and risk monitors cover the case
where labels arrive eventually~\cite{timans2025risk,zhang2026pprm,pouget2025suitability}.
\emph{Strength:} they target accuracy, the quantity operators care about, rather than
a distance between distributions. \emph{Limitation:} they are usually scored offline
on a fixed target set, and the time at which labels are released is rarely modelled.

\subsection{Delayed Labels and Adaptation}
Evaluating stream learners under delayed labels changes their measured
performance~\cite{grzenda2020delayed}, and some changes are undetectable until labels
arrive~\cite{zliobaite2010change}. PRODEM targets degradation detection with delayed
labels in fraud prediction~\cite{sethi2026prodem}. Proactive Hoeffding trees anticipate
incremental drift so that the model adapts before it degrades~\cite{guerrerocano2026proactive}.
Adaptation surveys organise reactive and proactive strategies~\cite{gama2014survey,lu2019learning}.
\emph{Strength:} they measure accuracy directly over time. \emph{Limitation:}
adaptation studies typically assume immediate (prequential) labels, and they rarely
compare update policies at matched label and refit budgets.

\subsection{Visual Analytics for Drift}
DriftVis~\cite{yang2020driftvis} and ConceptExplorer~\cite{wang2020conceptexplorer}
support interactive diagnosis of concept drift. Coordinated multiple views are an
established design for linking evidence across panels~\cite{roberts2007cmv}, and visual
analytics for deep learning has been surveyed~\cite{hohman2019visual}. Drift
visualization for specific learners has also been proposed~\cite{galmeanu2024svmviz}.
\emph{Strength:} they make drift legible to analysts. \emph{Limitation:} they are
evaluated mainly through case studies, and their displays rarely encode which outcomes
were still unknown at the time of a decision.

\subsection{Positioning}
Table~\ref{tab:related} summarises representative studies along the dimensions that
matter here. \idm{} does not propose a new detector. Its novelty lies in the
\emph{combination}: label-free input evidence and calibrated impact evidence are scored
against verified harm, under an enforced release delay, on streams where the
conditional law changes without any input change.

\begin{table*}[t]
\centering
\caption{Representative related studies organised by theme. \checkmark: addressed; --: not addressed in the evaluated setting.}
\label{tab:related}
\setlength{\tabcolsep}{3pt}
\footnotesize
\begin{tabular}{L{3.4cm}L{2.3cm}ccccL{4.3cm}}
\toprule
Study & Theme & Label-free & Delay modelled & Harm-linked & Explanation / view & Limitation for delayed-label harm monitoring \\
\midrule
Rabanser et al.~\cite{rabanser2019failing} & shift detection & \checkmark & -- & -- & -- & alarms not scored against accuracy loss \\
Greco et al.\ (DriftLens)~\cite{greco2025driftlens} & embedding drift & \checkmark & -- & -- & \checkmark & deep unstructured data; label-conditioned baselines \\
Lukats et al.~\cite{lukats2025benchmark} & detector benchmark & \checkmark & -- & -- & -- & real streams lack onset and harm truth \\
Sz\H{u}cs \& N\'emeth (DDE)~\cite{szucs2025dde} & drift localization & \checkmark & -- & -- & \checkmark & needs co-drifting correlated features \\
Hinder et al.~\cite{hinder2024surveyB} & drift explanation & \checkmark & -- & -- & \checkmark & descriptive, not causal \\
Garg et al.\ (ATC)~\cite{garg2022atc} & accuracy estimation & \checkmark & -- & \checkmark & -- & offline target sets; calibration-dependent \\
Kivim\"aki et al.~\cite{kivimaki2025confidence} & accuracy estimation & \checkmark & -- & \checkmark & -- & assumes calibration transfers under shift \\
Grzenda et al.~\cite{grzenda2020delayed} & delayed evaluation & -- & \checkmark & \checkmark & -- & evaluation protocol, not a monitor \\
Guerrero Cano et al.~\cite{guerrerocano2026proactive} & proactive adaptation & -- & -- & \checkmark & -- & prequential labels; incremental drift \\
Yang et al.\ (DriftVis)~\cite{yang2020driftvis} & visual analytics & \checkmark & -- & -- & \checkmark & case-study evaluation \\
\midrule
\textbf{\idm{} (this work)} & monitoring + evaluation & \checkmark & \checkmark & \checkmark & \checkmark & synthetic-heavy; no analyst study \\
\bottomrule
\end{tabular}
\end{table*}
